\chapter{Introduzione a SQL e Definizione dei Metadati}

\section{Nozioni introduttive}

Conclusa la fase di progettazione logica, il progettista ha tra le mani uno schema relazionale formale e matematicamente inattaccabile. Tuttavia, questo schema risiede ancora su carta. Il passo successivo consiste nell'istanziare fisicamente questo modello all'interno di un \textbf{DBMS} (Database Management System, come PostgreSQL, MySQL o Oracle).

Per comunicare con il motore del database si utilizza \textbf{SQL (Structured Query Language)}, il linguaggio standard universale per la gestione dei database relazionali.

\begin{dettagli}[Lo Standard SQL e la Portabilità]
Sebbene SQL sia il linguaggio standard per i database relazionali (codificato dagli enti internazionali ANSI e ISO), nella realtà commerciale ogni motore DBMS (PostgreSQL, MySQL, Oracle, ecc.) implementa un proprio "dialetto", introducendo estensioni e funzioni proprietarie non previste dallo standard. 

Per questo motivo, il manuale ufficiale di un buon DBMS inizia solitamente dichiarando esplicitamente il proprio livello di aderenza allo standard internazionale. Dal punto di vista dell'ingegneria del software, scrivere codice il più aderente possibile allo standard puro è una \textit{best practice} fondamentale: assicura un'elevata \textbf{portabilità} del sistema, garantendo la possibilità di migrare l'infrastruttura da un DBMS a un altro con sforzi e costi di riscrittura minimi. Nel corso della progettazione si farà riferimento esclusivamente alle direttive dell'SQL standard.
\end{dettagli}

\subsection*{Il concetto di Metadato e il Catalogo di Sistema}

Prima di poter inserire i dati reali dell'applicazione (es. Mario Rossi, matricola 12345), il DBMS deve conoscere la "forma" che questi dati avranno, le regole a cui devono sottostare e i legami che li uniscono. 
Queste informazioni strutturali prendono il nome di \textbf{Metadati} (letteralmente: \textit{dati che descrivono altri dati}).

Quando si definiscono gli schemi e i vincoli in SQL, il motore del database non sta ancora salvando informazioni applicative, ma sta popolando il proprio \textbf{Catalogo di Sistema} (o \textit{Data Dictionary}). Il catalogo è un set di tabelle interne, invisibili all'utente finale, che il DBMS usa per ricordare quante tabelle esistono, come si chiamano le colonne, quali sono i tipi di dato ammessi e quali vincoli di chiave primaria o esterna devono essere imposti.

\subsection*{L'Anatomia di SQL: Un Linguaggio Dichiarativo}

A differenza di linguaggi imperativi come Java, C o Python, in cui il programmatore deve descrivere \textit{passo dopo passo} come eseguire un'operazione, SQL è un linguaggio \textbf{dichiarativo}: il progettista si limita a dichiarare \textit{cosa} vuole ottenere o \textit{come} deve essere fatta la struttura; sarà il motore di ottimizzazione interno del DBMS a decidere la strategia algoritmica migliore per allocare la memoria o estrarre i dati.

Per gestire l'intero ciclo di vita del dato, il linguaggio SQL è suddiviso in tre grandi sotto-insiemi funzionali (sub-languages):

\begin{enumerate}
    \item \textbf{DDL (Data Definition Language):}
    È il linguaggio di definizione dei metadati. Viene utilizzato esclusivamente per manipolare la struttura del database (il contenitore). I suoi comandi principali sono:
    \begin{itemize}
        \item \texttt{CREATE}: per generare nuove tabelle, indici o viste.
        \item \texttt{ALTER}: per modificare la struttura di una tabella esistente (es. aggiungere una colonna).
        \item \texttt{DROP}: per distruggere fisicamente una struttura e i suoi dati.
    \end{itemize}

    \item \textbf{DML (Data Manipulation Language):}
    È il linguaggio utilizzato per manipolare il contenuto (i dati veri e propri) all'interno delle strutture precedentemente create col DDL. I comandi principali sono:
    \begin{itemize}
        \item \texttt{INSERT}: per inserire nuove tuple.
        \item \texttt{UPDATE}: per aggiornare i valori di tuple esistenti.
        \item \texttt{DELETE}: per rimuovere tuple.
    \end{itemize}

    \item \textbf{DQL (Data Query Language):}
    Spesso considerato parte del DML, è l'istruzione \texttt{SELECT}, il cuore pulsante dei database relazionali, utilizzata per interrogare il motore e filtrare, aggregare o unire (\texttt{JOIN}) i dati.
\end{enumerate}

\section{I Domini di Base (Tipi di Dato) nello Standard SQL}

Nella creazione di una tabella tramite DDL, ogni attributo deve essere obbligatoriamente associato a un \textbf{Dominio} (tipo di dato). Il dominio definisce lo spazio dei valori ammissibili per quella colonna, lo spazio fisico occupato su disco e le operazioni matematiche o logiche consentite.

Lo Standard SQL classifica i tipi di dato in grandi famiglie. Di seguito vengono presentati i domini fondamentali e maggiormente utilizzati.

\subsubsection*{Dati di tipo Stringa (Caratteri)}
Servono per memorizzare testo, codici alfanumerici o numeri su cui non si devono effettuare operazioni matematiche (es. numeri di telefono, CAP, codici fiscali).

\begin{itemize}
    \item \texttt{\textbf{CHAR(n)}}: Stringa a lunghezza \textbf{fissa} pari a $n$ caratteri. Se si inserisce una stringa più corta, il DBMS la "riempie" (padding) con spazi vuoti fino a raggiungere $n$.
    \begin{itemize}
        \item \textit{Uso ingegneristico:} Estremamente veloce per la ricerca. Si usa \textbf{solo} quando il dato ha lunghezza rigidamente costante (es. \texttt{CHAR(16)} per il Codice Fiscale, \texttt{CHAR(2)} per la sigla della Provincia).
    \end{itemize}
    
    \item \texttt{\textbf{VARCHAR(n)}}: Stringa a lunghezza \textbf{variabile} con limite massimo di $n$ caratteri. Se si inserisce una stringa più corta, occupa su disco solo lo spazio effettivo più un marcatore di fine stringa.
    \begin{itemize}
        \item \textit{Uso ingegneristico:} Lo standard \textit{de facto} per quasi tutto il testo (\texttt{VARCHAR(50)} per il nome, \texttt{VARCHAR(255)} per l'email, ad esempio). Fa risparmiare enorme spazio su disco, a costo di una microscopica latenza in più durante l'allocazione della memoria.
    \end{itemize}
\end{itemize}

\subsubsection*{Dati di tipo Numerico}
SQL distingue rigorosamente tra numeri esatti (dove 1.1 + 1.1 fa esattamente 2.2) e numeri approssimati (soggetti a virgola mobile e possibili errori di precisione).

\begin{itemize}
    \item \textbf{Numerici Esatti Interi:}
    \begin{itemize}
        \item \texttt{\textbf{INT}} o \texttt{\textbf{INTEGER}}: Numero intero standard con segno (solitamente a 32 bit, range da circa -2 a +2 miliardi). È il re indiscusso delle chiavi primarie surrogate e dei contatori.
        \item \texttt{\textbf{SMALLINT}}: Numero intero "piccolo" (solitamente a 16 bit). Utile per risparmiare RAM se il range di valori è sicuramente ridotto (es. l'anno di immatricolazione, o valori da 1 a 100).
    \end{itemize}
    
    \item \textbf{Numerici Decimali Esatti:}
    \begin{itemize}
        \item \texttt{\textbf{NUMERIC(p, s)}} o \texttt{\textbf{DECIMAL(p, s)}}: Numeri a virgola fissa. Il parametro $p$ (\textit{precision}) indica il numero totale di cifre, mentre $s$ (\textit{scale}) indica quante di queste cifre stanno dopo la virgola. 
        \item \textit{Esempio:} \texttt{NUMERIC(5, 2)} può memorizzare valori da -999.99 a 999.99.
        \item \textit{Uso ingegneristico:} Obbligatorio per i dati \textbf{finanziari e monetari} (stipendi, prezzi, bilanci). Garantisce che non ci siano mai arrotondamenti imprevisti.
    \end{itemize}
    
    \item \textbf{Numerici Approssimati:}
    \begin{itemize}
        \item \texttt{\textbf{REAL}}, \texttt{\textbf{DOUBLE PRECISION}}, \texttt{\textbf{FLOAT}}: Numeri a virgola mobile (standard IEEE 754).
        \item \textit{Uso ingegneristico:} Si usano \textbf{solo} per calcoli scientifici, coordinate GPS o misurazioni fisiche dove la scala è immensa e una micro-approssimazione (es. 1.9999999 al posto di 2) è tollerabile. Da evitare assolutamente per la valuta.
    \end{itemize}
\end{itemize}

\subsubsection*{Dati di tipo Temporale}
\begin{itemize}
    \item \texttt{\textbf{DATE}}: Memorizza solo la data (Anno, Mese, Giorno). Formato: \texttt{'YYYY-MM-GG'}. Ottimo per date di nascita o date di assunzione.
    \item \texttt{\textbf{TIME}}: Memorizza solo l'orario (Ore, Minuti, Secondi). Formato: \texttt{'HH-MI-SS'}.
    \item \texttt{\textbf{TIMESTAMP}} o \texttt{\textbf{DATETIME}}: Memorizza data e ora esatte in un unico campo. Formato: \texttt{'YYYY-MM-GG HH-MI-SS'}. Fondamentale per i log di sistema, la data di creazione di un record o per tracciare transazioni.
\end{itemize}

\begin{dettagli}[La Rappresentazione Fisica dei Tipi Temporali]
Nonostante in fase di interrogazione e inserimento le date vengano espresse in formato testuale (es. \texttt{'YYYY-MM-DD'}), a livello di immagazzinamento fisico i tipi temporali \textbf{non sono stringhe}. 

Il DBMS analizza la stringa immessa e la converte in un numero (spesso un intero a 32 o 64 bit che rappresenta i giorni o i millisecondi trascorsi da una specifica data di riferimento, detta \textit{Epoch}). 
Questa astrazione matematica è cruciale per tre motivi:
\begin{enumerate}
    \item \textbf{Performance aritmetica:} Permette al DBMS di sommare giorni o calcolare intervalli tramite velocissime operazioni matematiche a livello di CPU, evitando il costoso \textit{parsing} testuale.
    \item \textbf{Ottimizzazione spaziale:} Un \texttt{DATE} occupa tipicamente 4 byte, contro i 10 byte di un \texttt{CHAR(10)}.
    \item \textbf{Validazione semantica:} Il DBMS controlla nativamente l'esistenza reale della data (es. impedendo l'inserimento del 30 Febbraio).
\end{enumerate}
\end{dettagli}

\subsubsection*{Dati di tipo Booleano}
\begin{itemize}
    \item \texttt{\textbf{BOOLEAN}}: Ammette solo i valori di verità \texttt{TRUE} o \texttt{FALSE} (più il valore \texttt{NULL} se il campo non è obbligatorio). Ottimo per i flag (es. \texttt{is\_attivo}, \texttt{is\_pagato}).
\end{itemize}

\begin{warningbox}[Il NULL non è un tipo di dato]
È fondamentale ricordare che \texttt{NULL} non appartiene a nessun dominio specifico, ma è uno \textbf{stato} universale. \texttt{NULL} in SQL non significa "zero" né "stringa vuota", ma rappresenta "informazione mancante, sconosciuta o inapplicabile"; intuitivamente, si può pensare come una variabile incognita (come in una equazione matematica). Qualsiasi tipo di dato descritto sopra, in assenza di vincoli, può assumere lo stato \texttt{NULL}.
\end{warningbox}

\begin{dettagli}[Il supporto al tipo BOOLEAN e l'approccio legacy]
Sebbene lo standard SQL preveda nativamente il dominio \texttt{BOOLEAN}, nella pratica ingegneristica è necessario tenere conto delle limitazioni specifiche di alcuni DBMS. Il caso storico più celebre è Oracle Database, che non ha supportato il tipo booleano a livello di tabella fino alla versione 23c (2023).

Quando si opera su sistemi legacy o DBMS privi di questo dominio, la \textit{best practice} architetturale consiste nel "simulare" il booleano utilizzando un tipo di dato minimale, come \texttt{CHAR(1)} o \texttt{SMALLINT}/\texttt{NUMBER(1)}, e "blindarlo" tramite un vincolo di colonna personalizzato (il vincolo \texttt{CHECK}, che vedremo a breve) che limiti esplicitamente l'inserimento a due soli valori convenzionali.
\end{dettagli}

\section{Operatori Logici e di Confronto}

Per imporre regole di business (tramite vincoli \texttt{CHECK}) o per filtrare i dati durante le interrogazioni, SQL utilizza espressioni condizionali. Un'espressione condizionale valuta il contenuto di una cella e restituisce un risultato booleano a tre vie: \texttt{TRUE}, \texttt{FALSE} oppure \texttt{UNKNOWN} (se coinvolge un valore \texttt{NULL}).

Di seguito i principali operatori standard previsti dal linguaggio:

\subsection*{Operatori di Confronto Standard}
Si applicano a stringhe, numeri e date:
\begin{itemize}
    \item \texttt{=} (Uguale)
    \item \texttt{<>} oppure \texttt{!=} (Diverso)
    \item \texttt{<}, \texttt{>}, \texttt{<=}, \texttt{>=} (Minore, Maggiore e relativi uguali)
\end{itemize}

\subsection*{Operatori Logici (Algebra di Boole)}
Permettono di combinare più condizioni elementari:
\begin{itemize}
    \item \texttt{AND}: Restituisce \texttt{TRUE} solo se \textit{tutte} le condizioni combinate sono vere.
    \item \texttt{OR}: Restituisce \texttt{TRUE} se \textit{almeno una} delle condizioni è vera.
    \item \texttt{NOT}: Inverte il risultato logico della condizione che lo segue.
\end{itemize}

\subsection*{Operatori Speciali di SQL}
Sono costrutti sintattici avanzati che semplificano enormemente la scrittura del codice, ottimizzando le performance del motore di ricerca:
\begin{itemize}
    \item \texttt{\textbf{BETWEEN} x \textbf{AND} y}: Verifica se un valore è compreso in un intervallo continuo (inclusi gli estremi \texttt{x} e \texttt{y}). Ottimo per le date o i range numerici.
    \begin{itemize}
        \item \textit{Esempio:} \texttt{voto BETWEEN 18 AND 30} equivale a \texttt{voto >= 18 AND voto <= 30}.
    \end{itemize}
    
    \item \texttt{\textbf{IN} (val1, val2, \dots)}: Verifica se un valore appartiene a una lista di elementi discreti.
    \begin{itemize}
        \item \textit{Esempio:} \texttt{provincia IN ('RM', 'MI', 'NA')} equivale a \texttt{provincia = 'RM' OR provincia = 'MI' OR provincia = 'NA'}.
    \end{itemize}
    
    \item \texttt{\textbf{LIKE} 'pattern'}: Permette il confronto testuale parziale (pattern matching) sulle stringhe. Utilizza due caratteri jolly (wildcard):
    \begin{itemize}
        \item \texttt{\%} (Percentuale): Sostituisce una sequenza di zero o più caratteri.
        \item \texttt{\_} (Underscore): Sostituisce esattamente un singolo carattere.
        \item \textit{Esempio:} \texttt{cognome LIKE 'R\%'} (Trova tutti i cognomi che iniziano per R); \texttt{nome LIKE '\_ario'} (Trova Mario, Dario, ma non Ilario).
    \end{itemize}
\end{itemize}

\begin{warningbox}[La trappola del NULL: IS NULL vs = NULL]
Per verificare se una cella è vuota, l'intuito del programmatore suggerirebbe di scrivere \texttt{colonna = NULL}. In SQL, questa sintassi è \textbf{errata} e restituirà sempre \texttt{UNKNOWN}, facendo fallire la condizione. Poiché \texttt{NULL} rappresenta un valore "incognito", è impossibile affermare matematicamente che un'incognita sia "uguale" a un'altra incognita. 

L'unico modo legittimo in SQL per verificare l'assenza o la presenza di un dato è utilizzare gli operatori specifici di stato:
\begin{itemize}
    \item \texttt{\textbf{IS NULL}} (Verifica se il campo è vuoto)
    \item \texttt{\textbf{IS NOT NULL}} (Verifica se il campo contiene un valore)
\end{itemize}
\end{warningbox}

\section{Istruzioni DDL}

Il Data Definition Language (DDL) permette di definire la struttura fisica del database. I due comandi fondamentali per istanziare l'architettura dei dati sono \texttt{CREATE SCHEMA} e \texttt{CREATE TABLE}.

\begin{bestpractice}[Convenzione Stilistica (Case Sensitivity)]
Lo standard SQL è \textit{case-insensitive} per quanto riguarda le parole chiave del linguaggio. Tuttavia, è convenzione universale nell'ingegneria del software scrivere le keyword in \textbf{MAIUSCOLO} (es. \texttt{CREATE TABLE}, \texttt{VARCHAR}) e gli identificatori definiti dall'utente in \textbf{minuscolo} (es. \texttt{studenti}, \texttt{matricola}). Questo garantisce la massima leggibilità del codice.
\end{bestpractice}

\subsection{Gli Schemi Logici: Architettura, Namespacing e Modularità}

Nel contesto dei sistemi di gestione di basi di dati relazionali (RDBMS), e in particolare in PostgreSQL, un database non è una collezione piatta di tabelle. Esiste un livello di astrazione intermedio fondamentale: lo \textbf{schema}.

Per comprendere l'architettura logica del database, il parallelismo più efficace è quello con il file system del sistema operativo Linux:
\begin{itemize}
    \item Il \textbf{Database} rappresenta l'intero disco rigido (o una partizione logica isolata).
    \item Lo \textbf{Schema} equivale a una \textbf{cartella (directory)} all'interno di quel disco.
    \item Le \textbf{Tabelle}, le viste, gli indici e i tipi personalizzati sono i singoli \textbf{file} memorizzati dentro quella specifica cartella.
\end{itemize}

Una tabella non può esistere nel vuoto cosmico del database; deve risiedere tassativamente all'interno di uno schema. Se non viene configurato uno schema personalizzato durante la creazione delle strutture, il DBMS inserisce automaticamente le entità in uno schema di default denominato \texttt{public}.

\begin{lstlisting}[language=SQL, title=Sintassi di Creazione di uno Schema]
-- Creazione di un nuovo contenitore logico isolato
CREATE SCHEMA ecommerce;
\end{lstlisting}

Mantenere l'intera applicazione all'interno dello schema \texttt{public} è una scelta accettabile in contesti didattici. Nello sviluppo di sistemi enterprise e architetture software reali, tuttavia, la suddivisione del database in molteplici schemi risponde a precise necessità ingegneristiche:

\begin{enumerate}
    \item \textbf{Namespacing (Risoluzione dei conflitti):} Esattamente come in Java si utilizzano i \textit{package} per evitare conflitti tra classi aventi lo stesso nome, nel database gli schemi isolano gli identificatori. È possibile definire una tabella \texttt{products} nello schema \texttt{sales} (ottimizzata per il catalogo) e un'altra tabella \texttt{products} nello schema \texttt{inventory} (ottimizzata per la logistica di magazzino).
    \item \textbf{Isolamento dei Privilegi (Sicurezza):} Dal punto di vista del backend, diversi moduli o microservizi si connettono al medesimo database tramite utenti (\textit{database roles}) differenti. È possibile configurare i permessi affinché l'utente del servizio di fatturazione abbia privilegi di lettura/scrittura esclusivamente nello schema \texttt{billing}, rimanendo completamente cieco e privo di accessi nei confronti dello schema \texttt{human\_resources}.
    \item \textbf{Modularità del Dominio (Monoliti Modulari):} Permette di mappare i confini logici del software (\textit{Bounded Contexts} del Domain-Driven Design) direttamente a livello di persistenza, mantenendo separati i moduli di business (es. uno schema \texttt{auth} per le credenziali, uno schema \texttt{tickets} per l'assistenza, uno schema \texttt{core} per le logiche principali).
    \item \textbf{Architetture Multi-Tenant (SaaS):} Nello sviluppo di software Cloud in cui più aziende clienti (\textit{tenant}) condividono la stessa infrastruttura, lo schema rappresenta una soluzione ottimale per l'isolamento dei dati. Per evitare il consumo asimmetrico di risorse dovuto all'accensione di database multipli, si utilizza un unico database centralizzato in cui ogni cliente possiede il proprio schema dedicato (\texttt{tenant\_a.orders}, \texttt{tenant\_b.orders}).
\end{enumerate}

\paragraph{Risoluzione dei nomi e il \texttt{search\_path}}

Quando il backend esegue una query priva di riferimenti espliciti come \texttt{SELECT * FROM users;}, il motore di PostgreSQL deve determinare in quale "cartella" cercare la tabella. Per fare questo, consulta una variabile di configurazione di sessione denominata \textbf{\texttt{search\_path}}, il cui comportamento ricalca fedelmente la variabile d'ambiente \texttt{PATH} dei sistemi operativi.

Di default, il percorso è impostato per scansionare prima lo schema omonimo dell'utente connesso e, subito dopo, lo schema \texttt{public}. Se una tabella si trova in uno schema non incluso nel \texttt{search\_path}, il codice applicativo è obbligato a invocare l'oggetto tramite il suo \textbf{nome qualificato completo} (\textit{Fully Qualified Name}), anteponendo il nome dello schema separato da un punto:

\begin{lstlisting}[language=SQL, title=Utilizzo dei Nomi Qualificati]
-- Accesso esplicito bypassando il search_path di default
SELECT * FROM billing.invoices WHERE id = 10540;
\end{lstlisting}

\begin{tip}[Modifica dinamica del percorso di ricerca]
È possibile alterare il percorso di ricerca per la sessione corrente tramite il comando \texttt{SET}. Questo evita di dover digitare i nomi qualificati per ogni singola query:
\texttt{SET search\_path TO ecommerce, public;}
\end{tip}

\begin{dettagli}[Ispezione degli Schemi tramite SQL standard]
Tutti i metadati relativi alla struttura logica del database sono memorizzati nei dizionari di sistema. Per ricavare l'elenco completo degli schemi configurati senza dipendere dai meta-comandi del client \texttt{psql}, si può interrogare la vista standard:
\begin{verbatim}
SELECT schema_name FROM information_schema.schemata;
\end{verbatim}
\end{dettagli}

\subsection{Creazione di una tabella}
Il comando \texttt{CREATE TABLE} istanzia fisicamente uno schema relazionale $S(A_1, \dots, A_k)$ associando a ciascun attributo $A_i$ il proprio dominio $dom(A_i)$.
Il nome assegnato alla tabella deve essere \textbf{unico} all'interno dello schema (o del database, se non si usano schemi espliciti).

La struttura del comando prevede una coppia di parentesi tonde all'interno delle quali si susseguono istruzioni separate da virgole. 

\begin{example}[Creazione tabella senza vincoli espliciti]
Di seguito si riporta la traduzione in SQL dello schema relazionale relativo agli studenti. In questa prima fase esplorativa, ci limitiamo a definire i nomi degli attributi e i rispettivi domini (tipi di dato), omettendo temporaneamente i vincoli (come la chiave primaria).

\begin{lstlisting}[language=SQL, basicstyle=\ttfamily\small, keywordstyle=\color{blue}\bfseries]
CREATE TABLE studenti (
    matricola CHAR(9),
    cf        CHAR(16),
    nome      VARCHAR(20),
    cognome   VARCHAR(20),
    dataN     DATE
);
\end{lstlisting}

Logicamente, il corpo del comando si divide in due blocchi sequenziali:

\begin{enumerate}
    \item \textbf{Definizione degli Attributi:} In questa prima fase si elencano le colonne. Per ogni colonna si deve specificare il nome, il tipo di dato e, opzionalmente, dei vincoli specifici applicabili alla singola cella (vincoli intra-attributo).
    \item \textbf{Definizione dei Vincoli di Tabella:} In questa seconda fase (che vedremo subito a seguire), si esplicitano i vincoli complessi che coinvolgono più colonne contemporaneamente (come chiavi primarie composte o chiavi esterne).
\end{enumerate}

\medskip

La sintassi per la definizione di un singolo attributo (la riga base) è la seguente:
$$ \texttt{<nome\_attributo> <tipo\_attributo> [vincoli\_attributo]} $$


Si noti come la scelta dei domini rispecchi quanto discusso nel paragrafo precedente:
\begin{itemize}
    \item \texttt{matricola} e \texttt{cf} utilizzano \texttt{CHAR} poiché hanno una lunghezza fissa e rigidamente definita (rispettivamente 9 e 16 caratteri).
    \item \texttt{nome} e \texttt{cognome} utilizzano \texttt{VARCHAR} per ottimizzare lo spazio su disco, dato che la lunghezza reale del testo inserito è variabile.
    \item \texttt{dataN} (data di nascita) utilizza il tipo nativo \texttt{DATE} per garantire validazione e correttezza semantica.
\end{itemize}
\end{example}

\subsection{I Vincoli Intra-attributo (o in-line)}

I vincoli in-line sono regole di integrità applicate direttamente alla definizione di un singolo attributo, sulla stessa riga in cui viene dichiarato il tipo di dato. Servono a limitare i valori ammissibili nella singola cella, agendo come prima barriera contro la corruzione dei dati.

Lo Standard SQL prevede quattro vincoli principali in-line:

\paragraph{NOT NULL (Obbligatorietà)}
Indica che il campo non può mai assumere lo stato \texttt{NULL}. All'atto dell'inserimento di una nuova riga, l'utente è obbligato a fornire un valore valido per questa colonna.
\begin{itemize}
    \item \textit{Uso:} Si applica a tutti i dati vitali che definiscono l'entità (es. il cognome di una persona, il titolo di un libro).
\end{itemize}

\paragraph{UNIQUE (Univocità di colonna)}
Impone che tutti i valori presenti in quella colonna siano distinti tra loro. Il DBMS rifiuterà qualsiasi tentativo di inserire una riga con un valore già esistente in un'altra riga per quel campo.
\begin{itemize}
    \item \textit{Uso:} Si applica alle chiavi naturali candidate (es. il Codice Fiscale, l'Email). Nello standard SQL, a meno di specifiche implementazioni dei singoli DBMS, una colonna \texttt{UNIQUE} consente l'inserimento di molteplici valori \texttt{NULL}, poiché \texttt{NULL} rappresenta una incognita e due incognite non sono considerate uguali tra loro.
\end{itemize}

\paragraph{PRIMARY KEY (Chiave Primaria a livello di colonna)}
Elegge l'attributo a identificatore principale della tabella. Specificare \texttt{PRIMARY KEY} su una colonna implica contemporaneamente i vincoli \texttt{NOT NULL} e \texttt{UNIQUE}. Una tabella può avere \textbf{una sola} chiave primaria.
\begin{itemize}
    \item \textit{Uso:} Si usa quando la chiave primaria è composta da un solo attributo (es. una chiave surrogata come \texttt{id} o una chiave naturale singola come \texttt{matricola}).
\end{itemize}

\paragraph{DEFAULT (Valore predefinito)}
Non è un vincolo di integrità in senso stretto, ma una direttiva di assegnazione. Se durante un inserimento (\texttt{INSERT}) l'utente omette il valore per questa colonna, il DBMS vi inserisce automaticamente il valore specificato dopo la keyword \texttt{DEFAULT}.
\begin{itemize}
    \item \textit{Uso:} Ottimo per impostare stati iniziali (es. \texttt{is\_attivo BOOLEAN DEFAULT TRUE}) o timestamp di creazione (es. \texttt{data\_inserimento DATE DEFAULT CURRENT\_DATE}).
\end{itemize}

\begin{example}[Evoluzione dello schema studenti con vincoli di colonna]
Riprendiamo lo schema della tabella \texttt{studenti} e applichiamo i vincoli intra-attributo per blindarne la coerenza semantica ed evitare anomalie:

\begin{lstlisting}[language=SQL, basicstyle=\ttfamily\small, keywordstyle=\color{blue}\bfseries]
CREATE TABLE studenti (
    matricola CHAR(9) PRIMARY KEY,
    cf        CHAR(16) UNIQUE NOT NULL,
    nome      VARCHAR(20) NOT NULL,
    cognome   VARCHAR(20) NOT NULL,
    dataN     DATE
);
\end{lstlisting}

\textbf{Analisi delle scelte effettuate:}
\begin{itemize}
    \item \texttt{matricola}: È l'identificatore principale. Diventa \texttt{PRIMARY KEY}, quindi sarà univoca e obbligatoria per definizione.
    \item \texttt{cf}: Il codice fiscale deve essere univoco (\texttt{UNIQUE}) per evitare duplicati nella realtà, ma deve anche essere obbligatorio (\texttt{NOT NULL}) per ogni studente registrato.
    \item \texttt{nome} e \texttt{cognome}: Viene applicato \texttt{NOT NULL} perché un'istanza di studente priva di identità anagrafica non avrebbe alcun significato logico.
    \item \texttt{dataN}: Rimane priva di vincoli espliciti; questo significa che il sistema accetterà il valore \texttt{NULL} qualora la data di nascita non fosse momentaneamente disponibile al momento della registrazione.
\end{itemize}
\end{example}

\subsection{I Vincoli di Tabella (Out-of-line)}

Fino a questo momento abbiamo visto i vincoli applicati direttamente sulla stessa riga della definizione dell'attributo (sintassi \textit{in-line}). Tuttavia, lo standard SQL permette (e in molti casi impone) di dichiarare i vincoli in una sezione dedicata alla fine del comando \texttt{CREATE TABLE}, separata dalla definizione delle colonne. Questa formulazione prende il nome di \textbf{Vincolo di Tabella} (o \textit{out-of-line}).

La sintassi generale per dichiarare un vincolo a livello di tabella è:
$$ \texttt{CONSTRAINT <nome\_vincolo> <TIPO\_VINCOLO> (colonna1, colonna2, \dots)} $$

Sebbene l'indicazione del nome del vincolo sia opzionale (se omessa, il DBMS genera internamente un nome alfanumerico casuale, es. \texttt{SYS\_C00142}), nella pratica si raccomanda di assegnare sempre un nome esplicito e semantico (es. \texttt{pk\_studenti} o \texttt{unq\_cf}). 

Nominare i vincoli è fondamentale per due operazioni architetturali:
\begin{enumerate}
    \item \textbf{Leggibilità degli errori:} Se un'applicazione viola una regola, il database genererà un'eccezione specificando il nome del vincolo. Leggere \textit{"Violazione del vincolo unq\_cf"} permette al programmatore di capire istantaneamente che c'è un problema di codici fiscali duplicati, senza dover fare debug alla cieca.
    \item \textbf{Manipolazione a posteriori (Enable/Disable):} Quando il vincolo viene dichiarato, il DBMS lo crea e lo attiva automaticamente. Tuttavia, conoscendone il nome, l'amministratore del database può manipolarlo in futuro senza distruggere i dati. Tramite istruzioni specifiche (es. \texttt{ALTER TABLE ... DISABLE CONSTRAINT}), un vincolo può essere momentaneamente spento. Questa operazione è comunissima durante l'importazione massiva di milioni di dati (\textit{bulk insert}), dove spegnere i controlli accelera drasticamente la scrittura, per poi riattivarli (\textit{ENABLE}) a importazione finita.
\end{enumerate}

La motivazione tecnica principale che rende indispensabile la sintassi a livello di tabella è la gestione dei vincoli multi-colonna. 

Nel nostro esempio precedente, la chiave primaria era formata da un solo attributo (\texttt{matricola}). Ma cosa accade se la progettazione logica produce una \textbf{chiave primaria composta} (es. una tabella che memorizza gli esami superati da uno studente in un determinato corso)?
In SQL è illegale scrivere \texttt{PRIMARY KEY} su due righe separate, perché il DBMS tenterebbe di creare \textit{due} chiavi primarie distinte per la stessa tabella. In questo scenario, \textbf{l'uso del vincolo di tabella diventa strutturalmente obbligatorio}, poiché permette di racchiudere più colonne sotto un'unica regola.

\begin{example}[Riscrittura e Chiavi Composte]
Mostriamo la sintassi applicata allo schema \texttt{studenti} e a una nuova tabella associativa \texttt{esami\_sostenuti}.

\begin{lstlisting}[language=SQL, basicstyle=\ttfamily\small, keywordstyle=\color{blue}\bfseries]
-- Esempio 1: Riscrittura con vincoli di tabella nominati
CREATE TABLE studenti (
    matricola CHAR(9),
    cf        CHAR(16) NOT NULL, -- NOT NULL resta solitamente in-line
    nome      VARCHAR(20) NOT NULL,
    cognome   VARCHAR(20) NOT NULL,
    dataN     DATE,
    
    -- Sezione Vincoli di Tabella
    CONSTRAINT pk_studenti PRIMARY KEY (matricola),
    CONSTRAINT uq_studenti_cf UNIQUE (cf)
);

-- Esempio 2: Obbligatorieta' del vincolo per PK composta
CREATE TABLE esame_sostenuto (
    studente_matr CHAR(9),
    codice_corso  CHAR(5),
    voto          INT NOT NULL,
    
    -- La PK unisce due attributi e deve essere definita qui
    CONSTRAINT pk_esami PRIMARY KEY (studente_matr, codice_corso)
);
\end{lstlisting}
\end{example}

\begin{bestpractice}[Convenzione e Best Practice]
Come regola pratica di sviluppo: il vincolo d'obbligatorietà (\texttt{NOT NULL}) si definisce \textbf{sempre e solo in-line} accanto all'attributo. Al contrario, i vincoli strutturali d'identità (\texttt{PRIMARY KEY}, \texttt{UNIQUE}) e quelli di navigabilità (\texttt{FOREIGN KEY}, che vedremo a breve) vengono quasi sempre traslati a fine tabella e nominati esplicitamente, per favorire l'ordine mentale e la manutenzione del codice.
\end{bestpractice}

\subsection{Classificazione Concettuale dei Vincoli}

Nei paragrafi precedenti abbiamo analizzato una classificazione prettamente \textit{sintattica} (vincoli \textit{in-line} vs \textit{out-of-line}), che riguarda esclusivamente la forma e il posizionamento del codice all'interno dell'istruzione \texttt{CREATE TABLE}. 

Esiste tuttavia una classificazione ben più importante, di natura \textbf{concettuale}, che categorizza i vincoli in base al loro "raggio d'azione", ovvero alla porzione di dati che il DBMS deve leggere per poter verificare se il vincolo è rispettato o violato. 
Secondo questa logica, i vincoli si dividono in quattro categorie a complessità crescente:

\begin{enumerate}
    \item \textbf{Vincoli di Dominio (o di colonna):} 
    Esprimono condizioni sui valori ammissibili per un singolo attributo, valutando la singola cella isolata dal resto. Sono i vincoli più leggeri da verificare per il motore del database.
    \begin{itemize}
        \item \textit{Osservazione sintattica:} Questo è tipicamente l'unico tipo di vincolo che ha senso logico definire \textit{in-line}.
        \item \textit{Esempi già incontrati:} Il vincolo di obbligatorietà (\texttt{NOT NULL}) e l'assegnazione di un valore predefinito (\texttt{DEFAULT}). Per verificare se un campo è nullo, il DBMS deve guardare solo quella specifica cella.
    \end{itemize}

    \item \textbf{Vincoli di tupla (o di riga):}
    Esprimono condizioni logiche che coinvolgono contemporaneamente più attributi appartenenti alla stessa riga (tupla). Per verificare il vincolo, il DBMS deve estrarre l'intera riga, ma non ha bisogno di confrontarla con le altre righe della tabella.

    \item \textbf{Vincoli Intra-relazionali (o di tabella):}
    Esprimono condizioni che coinvolgono più record all'interno della stessa tabella. Per validare il nuovo dato, il DBMS è costretto a scansionare o interrogare le altre righe già presenti.
    \begin{itemize}
        \item \textit{Esempi già incontrati:} \texttt{UNIQUE} e \texttt{PRIMARY KEY}. Per garantire che un Codice Fiscale sia univoco, non basta guardare la riga appena inserita; il database deve verificare che quel valore non sia già presente in nessun altro record della tabella.
    \end{itemize}

    \item \textbf{Vincoli Inter-relazionali (o di Schema):}
    Esprimono condizioni che legano tra loro record appartenenti a due o più tabelle distinte. Sono i vincoli più complessi e costosi computazionalmente, in quanto costringono il DBMS a "navigare" tra relazioni diverse per validare l'integrità strutturale del database.
\end{enumerate}

\begin{example}[Istanziazione della Tabella Esami e Classificazione dei Vincoli]
Esemplifichiamo la classificazione concettuale dei vincoli modellando la tabella \texttt{esami}, la quale memorizza la registrazione definitiva degli esami superati dagli studenti. 

Dal punto di vista della \textit{business logic}:
\begin{itemize}
    \item La chiave primaria deve essere composta dalla coppia \texttt{(matr, corso)}, garantendo che uno studente possa registrare un determinato esame una sola volta.
    \item Il voto deve essere unicamente compreso tra 18 e 30 (esame superato).
    \item La lode (valori 'S'/'N', default 'N') può essere assegnata solo se il voto è pari a 30.
    \item La matricola deve fare riferimento a uno studente reale.
\end{itemize}

Ecco il codice SQL DDL definitivo:

\begin{lstlisting}[language=SQL, basicstyle=\ttfamily\small, keywordstyle=\color{blue}\bfseries]
CREATE TABLE esame (
    matr  CHAR(9),
    corso VARCHAR(20) NOT NULL,
    data  DATE NOT NULL,
    voto  INTEGER NOT NULL,
    lode  CHAR(1) DEFAULT 'N',
    
    -- Sezione Vincoli di Tabella
    CONSTRAINT pk_esame PRIMARY KEY (matr, corso),
    
    CONSTRAINT chk_voto_superato CHECK (voto BETWEEN 18 AND 30),
    
    CONSTRAINT chk_lode_coerenza CHECK (lode IN ('S', 'N') AND (lode = 'N' OR voto = 30)),
    
    CONSTRAINT fk_esame_studente FOREIGN KEY (matr) 
        REFERENCES studente(matricola)
);
\end{lstlisting}



Analizziamo ogni singolo vincolo applicato alla tabella \texttt{esami}, esplicitando la categoria concettuale di appartenenza in base al suo spettro d'azione:

\begin{enumerate}
    \item \textbf{Vincoli di Dominio (o di colonna):}
    \begin{itemize}
        \item \texttt{NOT NULL} (su \texttt{corso}, \texttt{data}, \texttt{voto}): Sono vincoli di dominio in-line. Impongono che le singole celle non siano vuote. Il DBMS verifica la singola cella isolata.
        \item \texttt{DEFAULT 'N'} (su \texttt{lode}): È una direttiva di dominio che interviene sulla singola cella in assenza di dato.
        \item \texttt{chk\_voto\_superato} (\texttt{CHECK (voto BETWEEN 18 AND 30)}): È un vincolo di dominio out-of-line. Anche se scritto in fondo, analizza \textbf{esclusivamente} i valori dell'attributo \texttt{voto}, isolato dal resto della riga.
    \end{itemize}

    \item \textbf{Vincoli di n-pla (o di tupla/riga):}
    \begin{itemize}
        \item \texttt{chk\_lode\_coerenza}: Questo vincolo impone che il carattere sia 'S' o 'N' e che, qualora sia 'S', il voto sia tassativamente 30. Dal momento che l'espressione logica mette in relazione \textbf{due attributi distinti dello stesso record} (\texttt{lode} e \texttt{voto}), si tratta di un puro vincolo di n-pla. Il DBMS deve leggere l'intera riga per validarlo.
    \end{itemize}

    \item \textbf{Vincoli Intra-relazionali (o di Tabella):}
    \begin{itemize}
        \item \texttt{pk\_esame} (\texttt{PRIMARY KEY (matr, corso)}): Definisce l'identificatore. Per garantire che la coppia matricola-corso sia unica, il DBMS non può limitarsi a guardare la riga corrente, ma deve scansionare l'intera tabella \texttt{esame} per escludere duplicati. Agisce quindi a livello di \textbf{più record della stessa tabella}.
    \end{itemize}

    \item \textbf{Vincoli Inter-relazionali (o di Schema):}
    \begin{itemize}
        \item \texttt{fk\_esame\_studente} (\texttt{FOREIGN KEY}): È il vincolo che garantisce l'integrità referenziale. Stabilisce un legame tra la tabella corrente (\texttt{esami}) e una tabella esterna (\texttt{studente}). Il database, per accettare la riga, è costretto a saltare fuori dalla tabella corrente e verificare la presenza della matricola nello schema dello studente. \textbf{Questo è l'unico tipo di vincolo inter-relazionale che si può inserire in una tabella}.
    \end{itemize}
\end{enumerate}
\end{example}

\subsection{Validazione Avanzata: Pattern Matching e Regex}

Finora abbiamo visto vincoli di dominio (\texttt{CHECK}) basati su confronti matematici semplici (es. \texttt{prezzo > 0}). Tuttavia, quando si modellano stringhe di testo che devono rispettare rigorosi standard legali o di business (come un Codice Fiscale, un'Email o una Partita IVA), la semplice verifica della lunghezza o l'uso dell'operatore \texttt{LIKE} risulta insufficiente.

L'operatore \texttt{LIKE '\%@\%.\%'} può verificare rozzamente se una stringa contiene una chiocciola e un punto, ma non può impedire l'inserimento di valori insensati come \texttt{"@@@..."}.

Per risolvere questo problema, i moderni DBMS relazionali implementano il supporto alle \textbf{Espressioni Regolari (Regular Expressions o Regex)}, ovvero sequenze di caratteri che definiscono un pattern di ricerca estremamente preciso.

\paragraph{L'operatore POSIX in PostgreSQL}
Mentre lo standard SQL prevede la clausola \texttt{SIMILAR TO}, nell'ingegneria backend su PostgreSQL è considerato standard architetturale l'utilizzo dell'operatore nativo POSIX \textbf{\texttt{\textasciitilde}} (tilde) all'interno dei vincoli \texttt{CHECK}.

\begin{itemize}
    \item \texttt{colonna \textasciitilde\ 'pattern'}: La stringa deve rispettare il pattern (Case-sensitive).
    \item \texttt{colonna \textasciitilde* 'pattern'}: La stringa deve rispettare il pattern (Case-insensitive).
    \item \texttt{colonna !\textasciitilde\ 'pattern'}: La stringa \textbf{non} deve rispettare il pattern.
\end{itemize}

\begin{example}[La Difesa in Profondità nel Dominio SafeDrive]
Si consideri lo schema per memorizzare i dati di un perito assicurativo e di un cliente. I campi come il CAP o la Partita IVA devono essere validati.
\begin{lstlisting}[language=SQL, basicstyle=\ttfamily\small, keywordstyle=\color{blue}\bfseries]
CREATE TABLE customer (
    id BIGINT GENERATED ALWAYS AS IDENTITY PRIMARY KEY,
    cap_code CHAR(5) NOT NULL,
    fiscal_code CHAR(16) UNIQUE NOT NULL,
    
    -- Il CAP deve essere composto ESATTAMENTE da 5 numeri
    CONSTRAINT chk_cap_format CHECK (cap_code ~ '^[0-9]{5}$'),
    
    -- Il C.F. deve essere ESATTAMENTE di 16 caratteri (lettere o numeri)
    CONSTRAINT chk_fiscal_code CHECK (fiscal_code ~ '^[A-Z0-9]{16}$')
);
\end{lstlisting}

\textbf{Anatomia del Pattern \texttt{'\textasciicircum[0-9]\{5\}\$'}:}
\begin{itemize}
    \item \textbf{\texttt{\textasciicircum}} : Ancora di inizio. Indica che il pattern deve partire dal primissimo carattere della stringa.
    \item \textbf{\texttt{[0-9]}} : Classe di caratteri. Ammette esclusivamente le cifre da 0 a 9.
    \item \textbf{\texttt{\{5\}}} : Quantificatore. Impone che la classe precedente (le cifre) si ripeta esattamente 5 volte.
    \item \textbf{\texttt{\$}} : Ancora di fine. Indica che la stringa deve terminare immediatamente dopo il quinto numero, impedendo l'inserimento di caratteri extra alla fine.
\end{itemize}
\end{example}

\begin{observation}[L'importanza architetturale (Defense in Depth)]
Applicare le regex nel database non rende superflua la validazione a livello applicativo (es. tramite annotazioni \texttt{@Pattern} in Java/Spring Boot). Il livello Java deve continuare a validare i dati per restituire errori HTTP chiari (es. 400 Bad Request) all'utente. 

Tuttavia, il vincolo \texttt{CHECK} sul database rappresenta la \textbf{difesa in profondità}: garantisce che i dati non vengano mai corrotti nemmeno se un amministratore effettua un inserimento manuale, se si esegue un'importazione massiva da un file CSV ignorando il backend, o se si verifica un bug nel codice dell'applicativo. Il database è l'ultima fonte di verità e deve sapersi difendere da solo.
\end{observation}

\subsection{Integrità Referenziale e Chiavi Esterne (FOREIGN KEY)}

Il vincolo di \textbf{Chiave Esterna (FOREIGN KEY)} è il più importante tra i vincoli \textbf{inter-relazionali} (o di schema). Esso stabilisce una relazione di dipendenza tra due tabelle, imponendo che i valori presenti in un set di attributi della tabella corrente (\textit{tabella figlia} o \textit{referenziante}) debbano obbligatoriamente esistere all'interno del set di attributi che costituisce la chiave primaria (o un'altra chiave univoca) di un'altra tabella (\textit{tabella padre} o \textit{referenziata}).

% Garantisce l'\textbf{Integrità Referenziale}: non possono esistere "puntatori orfani" nel database.

Come abbiamo visto nel paragrafo precendente, la sintassi standard a livello di tabella è la seguente:
\begin{lstlisting}[language=SQL]
CONSTRAINT <nome_fk> FOREIGN KEY (colonna_locale_1, ...)
    REFERENCES <tabella_padre> (colonna_padre_1, ...)
\end{lstlisting}

\begin{warningbox}[Ordine Posizionale degli Attributi Referenziati]
Nel caso in cui la relazione coinvolga chiavi composte da più attributi (es. una tabella figlia che punta alla nostra tabella \texttt{esami}, la cui Primary Key è la coppia \texttt{matr, corso}), l'\textbf{ordine} in cui le colonne vengono dichiarate all'interno delle parentesi tonde è tassativo.

L'associazione effettuata dal motore SQL \textbf{non avviene per nome, ma è strettamente posizionale} (uno-a-uno): il primo attributo elencato in \texttt{FOREIGN KEY} viene mappato sul primo attributo elencato in \texttt{REFERENCES}, il secondo sul secondo, e così via. Un'inversione accidentale dell'ordine genererà un errore di compilazione (se i domini non corrispondono) o, nel peggiore dei casi, una drammatica corruzione semantica dei dati referenziati.
\end{warningbox}

\subsubsection{Il problema delle modifiche sui dati referenziati}
Il vero fulcro ingegneristico della chiave esterna emerge quando si tenta di alterare lo stato della tabella padre. Supponiamo che nel nostro database esista uno studente con matricola \texttt{'12345'} e che nella tabella \texttt{esami} vi siano tre record associati a tale matricola. 

Cosa accade se un operatore tenta di \textbf{cancellare} lo studente \texttt{'12345'} dalla tabella \texttt{studenti}, o di \textbf{modificarne} la matricola? 
Se il DBMS permettesse questa operazione alla leggera, i tre record della tabella \texttt{esami} punterebbero a una matricola inesistente.

Per risolvere questo problema, lo Standard SQL impone al progettista di specificare esplicitamente una \textbf{politica di reazione} sia per la cancellazione (\texttt{ON DELETE}) sia per l'aggiornamento (\texttt{ON UPDATE}).

Il progettista può scegliere tra quattro comportamenti distinti, da dichiarare in coda al vincolo di chiave esterna:

\begin{enumerate}
    \item \textbf{RESTRICT / NO ACTION (Comportamento di Default):}
    Il DBMS blocca l'operazione alla radice. Se si tenta di cancellare lo studente padre, il database lancia un'eccezione di violazione di vincolo e rifiuta l'operazione, costringendo l'utente a cancellare manualmente prima tutti gli esami correlati.
    \begin{itemize}
        \item \texttt{Sintassi:} \texttt{ON DELETE RESTRICT}
    \end{itemize}

    \item \textbf{CASCADE (Cancellazione/Modifica a cascata):}
    L'operazione eseguita sul padre si ripercuote automaticamente sui figli. Se cancello lo studente, il DBMS cancella autonomamente tutte le righe dei suoi esami. Se modifico la matricola del padre, il DBMS aggiorna la matricola in tutti gli esami figli.
    \begin{itemize}
        \item \texttt{Sintassi:} \texttt{ON DELETE CASCADE}
    \end{itemize}

    \item \textbf{SET NULL (Annullamento del puntatore):}
    Se il padre viene eliminato o modificato, i record figli sopravvivono, ma il valore della loro chiave esterna viene impostato automaticamente a \texttt{NULL}. 
    \begin{itemize}
        \item \texttt{Sintassi:} \texttt{ON DELETE SET NULL}
        \item \textit{Vincolo strutturale:} Questa politica è applicabile \textbf{solo se} le colonne della chiave esterna nella tabella figlia ammettono il valore \texttt{NULL} (ovvero non sono marcate come \texttt{NOT NULL} o non fanno parte della Primary Key).
    \end{itemize}

    \item \textbf{SET DEFAULT (Ripristino del valore predefinito):}
    Simile a \texttt{SET NULL}, ma i campi della chiave esterna nei record figli assumono il valore specificato nella clausola \texttt{DEFAULT} di quella colonna.
    \begin{itemize}
        \item \texttt{Sintassi:} \texttt{ON DELETE SET DEFAULT}
    \end{itemize}
\end{enumerate}

\begin{observation}[Euristiche di scelta delle politiche]
La selezione della politica referenziale dipende esclusivamente dal significato semantico della relazione:
\begin{itemize}
    \item Si usa \textbf{\texttt{CASCADE}} quando la tabella figlia descrive entità deboli che non hanno senso di esistere senza il padre (es. le righe di un ordine d'acquisto: se cancello l'ordine, devo distruggere le sue righe).
    \item Si usa \textbf{\texttt{RESTRICT}} quando l'eliminazione del padre sarebbe un errore operativo grave (es. non devo poter cancellare un corso di laurea se ci sono ancora studenti iscritti ad esso).
    \item Si usa \textbf{\texttt{SET NULL}} quando l'entità figlia ha vita propria ma perde il legame con il padre (es. se una scuderia di Formula 1 fallisce e si cancella la scuderia, i piloti rimangono nel database ma con la colonna \texttt{id\_scuderia} impostata a \texttt{NULL}).
\end{itemize}
\end{observation}

\begin{tip}[L'obsolescenza di ON UPDATE nell'era delle Chiavi Surrogate]
Nei database legacy o nei testi accademici che fanno largo uso di Chiavi Naturali, è considerata \textit{best practice} definire comportamenti misti, come:
\begin{lstlisting}[language=SQL, basicstyle=\ttfamily\small, keywordstyle=\color{blue}\bfseries]
CONSTRAINT fk_esami_studenti FOREIGN KEY (matr) REFERENCES studenti(matricola)
    ON DELETE RESTRICT   -- Protegge dalla perdita accidentale dei dati
    ON UPDATE CASCADE    -- Propaga la correzione della matricola
\end{lstlisting}

Tuttavia, l'ingegneria del software moderna impone quasi universalmente l'uso di \textbf{Chiavi Surrogate} (es. un \texttt{id} numerico auto-incrementante) come Primary Key. 
La regola fondamentale di una chiave surrogata è la sua \textbf{immutabilità assoluta}: un \texttt{id} nasce e muore con il record e non ha alcun motivo di business per essere alterato.

Di conseguenza, il problema di "cosa fare quando la chiave del padre cambia" cessa semplicemente di esistere. Nell'architettura backend odierna (specialmente lavorando con ORM come Spring Data JPA), la clausola \texttt{ON UPDATE} è quasi sempre omessa, poiché definire una politica di aggiornamento su una colonna immutabile genererebbe unicamente \textbf{codice morto}. 
Il progettista moderno concentra l'attenzione esclusivamente sulla gestione del ciclo di vita dei dati (\texttt{ON DELETE}), affidandosi al \texttt{RESTRICT} predefinito per impedire matematicamente qualsiasi tentata, e illegittima, manipolazione degli identificatori di sistema.
\end{tip}

\begin{dettagli}[Design Pattern: Il Record Fittizio (Dummy Record)]
Durante la modellazione fisica, la gestione dei valori \texttt{NULL} nelle chiavi esterne può generare complessità in fase di interrogazione dell'applicativo (necessità di \textit{Outer Join} e gestione delle incognite). 

Per ovviare a questo problema, i progettisti adottano spesso il pattern del \textbf{Record Fittizio}. La tecnica consiste nell'inserire nella tabella padre una tupla convenzionale (es. con Primary Key \texttt{0} o \texttt{-1} e descrizione "Non Assegnato"). 
Nella tabella figlia, la colonna dedicata alla chiave esterna viene dichiarata \texttt{NOT NULL} (impedendo l'incognita) e le viene assegnato come \texttt{DEFAULT} l'identificatore del record fittizio.

\textit{Vantaggi architetturali:}
\begin{itemize}
    \item \textbf{Integrità totale:} Ogni riga figlia punta sempre a un record padre esistente, eliminando il concetto di "puntatore vuoto".
    \item \textbf{Sinergia referenziale:} Rende implementabile logicamente la direttiva \texttt{ON DELETE SET DEFAULT}. Se un'entità padre viene eliminata, i record figli non diventano \texttt{NULL}, ma passano al record "Sconosciuto".
\end{itemize}

\smallskip

\textbf{Attenzione Architetturale (OLTP vs OLAP):} \\
Questo pattern è considerato la \textit{prassi d'oro} nel mondo della \textbf{Business Intelligence e dei Data Warehouse (OLAP)}, in particolare nella modellazione dimensionale (Star Schema), dove le tabelle dei fatti non devono mai contenere chiavi esterne \texttt{NULL} per garantire la massima performance delle \textit{Inner Join} durante i calcoli di aggregazione.

Al contrario, nei moderni database transazionali \textbf{(OLTP)} operati tramite \textbf{ORM} (es. Hibernate/Spring Boot), questo approccio è considerato un \textbf{Anti-Pattern}. Negli OLTP l'assenza di relazione deve essere modellata in modo semanticamente puro (accettando il \texttt{NULL} fisico) ed è delegata alla logica applicativa (es. l'uso di \texttt{Optional} in Java), poiché l'inserimento di record fittizi "inquinerebbe" il dominio, obbligando a filtri costanti per nascondere i dati inesistenti agli utenti finali.
\end{dettagli}

\begin{warningbox}[]
    In alcuni contesti didattici, con \textbf{Vincolo di Integrità Referenziale} si intende che la FK nella tabella figlia debba \textbf{sempre} puntare ad una tupla nella tabella padre, eventualmentente un dummy record (ad esempio, nel caso in cui la FK sia \texttt{NULL}). 
\end{warningbox}

\subsection{Identificatori Univoci: Chiavi Surrogate e il costrutto IDENTITY}

Come sappiamo dai capitoli precedenti, per garantire un'identificazione robusta e immutabile nel tempo, l'ingegneria del software ricorre quasi sempre alle \textbf{Chiavi Surrogate}: attributi artificiali, tipicamente un numero intero progressivo (es. 1, 2, 3...), privi di qualsiasi significato logico per l'utente, ma utilizzati dal motore relazionale per collegare le tabelle tra loro in modo sicuro.

\begin{dettagli}[L'anti-pattern del MAX(id) + 1]
    Perché non possiamo semplicemente calcolare il nuovo ID con una query del tipo \textit{"Prendi il numero massimo attuale e aggiungi 1"} (\texttt{MAX(id) + 1})?
    In un'applicazione aziendale, centinaia di utenti potrebbero tentare di registrarsi contemporaneamente. Se due transazioni eseguissero \texttt{MAX(id) + 1} nello stesso esatto millisecondo, leggerebbero entrambe lo stesso numero (es. 100), calcolerebbero lo stesso nuovo ID (101) e cercherebbero di inserirlo. Il database andrebbe in crash per violazione della chiave primaria. 

    Per risolvere questo problema serviva un generatore di numeri \textbf{atomico e isolato} dalle transazioni: un oggetto capace di "staccare un biglietto" diverso per ogni richiesta, garantendo che due richieste simultanee non ricevano mai lo stesso numero. Questo oggetto matematico prende il nome di \textbf{Sequenza} (\texttt{SEQUENCE}).
\end{dettagli}

\paragraph{La Soluzione dello Standard SQL Moderno}

Nello standard SQL moderno (pienamente supportato da PostgreSQL 10+, Oracle 12c+ e DB2), non è più necessario creare e gestire manualmente l'oggetto Sequenza. Il linguaggio offre un costrutto dichiarativo estremamente elegante da inserire direttamente nella definizione della tabella: la clausola \textbf{\texttt{IDENTITY}}.

Quando si dichiara una colonna come \texttt{IDENTITY}, il DBMS crea in automatico, "sotto il cofano", una sequenza blindata e la associa indissolubilmente a quell'attributo.

\paragraph{Sintassi e Comportamento}

Esistono due varianti sintattiche principali che ne determinano il comportamento durante l'istruzione \texttt{INSERT}:

\textbf{1. GENERATED ALWAYS AS IDENTITY}
Questa è la forma più restrittiva e sicura. Il motore relazionale prende il controllo totale della colonna. 
\begin{lstlisting}[language=PostgreSQL, basicstyle=\ttfamily\small, keywordstyle=\color{blue}\bfseries]
CREATE TABLE utente (
    id_utente INT GENERATED ALWAYS AS IDENTITY PRIMARY KEY,
    nome VARCHAR(50),
    email VARCHAR(100)
);
\end{lstlisting}
In fase di popolamento, lo sviluppatore \textbf{deve obbligatoriamente ignorare} la colonna \texttt{id\_utente}:
\begin{lstlisting}[language=SQL, basicstyle=\ttfamily\small, keywordstyle=\color{blue}\bfseries]
-- L'ID verra' generato e inserito autonomamente dal motore
INSERT INTO utente (nome, email) VALUES ('Mario Rossi', 'mario@email.it');
\end{lstlisting}
Se si tenta di forzare l'inserimento manuale di un ID (es. \texttt{VALUES (15, 'Mario', ...)}), il DBMS bloccherà l'operazione restituendo un errore, preservando così l'integrità del contatore numerico.

\textbf{2. GENERATED BY DEFAULT AS IDENTITY}
Questa variante agisce come un paracadute. Il motore genererà il numero progressivo in automatico, ma se l'applicazione ha un motivo specifico per forzare l'inserimento di un ID manuale (ad esempio, durante il ripristino di un backup o la migrazione da un vecchio database), il DBMS lo accetterà.

\begin{dettagli}[I "buchi" di numerazione]
Una delle regole più controintuitive per chi si affaccia ai database riguarda i "buchi" (\textit{gaps}) nelle sequenze. 
Se una transazione "stacca" l'ID numero 105 per inserire un utente, ma l'inserimento fallisce (es. perché l'email violava un vincolo \texttt{UNIQUE}), o se la riga viene successivamente cancellata con una \texttt{DELETE}, \textbf{l'ID 105 è perso per sempre}. L'inserimento successivo riceverà il 106. 

Questo comportamento non è un difetto, ma una caratteristica fondamentale per garantire alte prestazioni. Se le sequenze dovessero aspettare l'esito di ogni transazione per riutilizzare i numeri scartati, creerebbero colli di bottiglia catastrofici. Le chiavi surrogate devono unicamente garantire l'univocità, non la continuità numerica assoluta.
\end{dettagli}